\documentclass[11pt,oneside]{book}

% ---------- Packages ----------
\usepackage[T1]{fontenc}
\usepackage[utf8]{inputenc}
\usepackage{lmodern}
\usepackage{amsmath,amssymb,amsthm} % math + theorems [web:6]
\usepackage{mathtools}
\usepackage{geometry}
\usepackage{hyperref}
\usepackage{enumitem}

\geometry{margin=1in}


% ============ hyperlinks and clever reference
\usepackage{hyperref} % hyperref must be placed before cleveref
\hypersetup{colorlinks=true, citecolor=blue, urlcolor=blue, linkcolor=blue}
\usepackage{cleveref}
  	\crefname{figure}{Figure}{Figures}
  	\crefname{table}{Table}{Tables}
  	\crefname{equation}{Eq.}{Eqs.}
  	\crefname{section}{Section}{Sections}
  	\crefname{subsection}{Section}{Sections}
  	\crefname{subsubsection}{Section}{Sections}
  	\crefname{algorithm}{Algorithm}{Algorithms}
    \crefname{appendix}{Appendix}{Appendices}

% ---------- Document ----------
\begin{document}

\frontmatter
\tableofcontents

\mainmatter

% ==============================
\chapter{Spin vs Kramers restricted MP2}
\label{ch:SR-vs-KR}

\section{Spin-unrestricted non-relativistic MP2 energy}

The non-relativistic MP2 energy in spin-orbitals, assuming Einstein summation, is
\begin{align}
    E_{\text{MP2}} &= \frac{1}{4} \bar{t}_{ab}^{ij} \bar{g}_{ij}^{ab} .
\end{align}
Here, $\bar{g}_{ij}^{ab}$ denotes the antisymmetrized electron repulsion integrals,
\begin{align}
    \bar{g}_{ij}^{ab} 
    &= g_{ij}^{ab} - g_{ji}^{ab} = \langle ij | ab \rangle - \langle ji | ab \rangle = \langle ij \Vert ab \rangle,\\
    g_{ij}^{ab} &=  \langle ij | ab \rangle = \iint \phi_i^*(\mathbf{r1}) \phi_j^*(\mathbf{r2}) \frac{1}{|\mathbf{r1} - \mathbf{r2}|} \phi_a(\mathbf{r1}) \phi_b(\mathbf{r2}) d\mathbf{r1} d\mathbf{r2}
\end{align}
and
\begin{align}
    \bar{t}_{ab}^{ij} = \frac{\bar{g}_{ab}^{ij}}{\epsilon_i + \epsilon_j - \epsilon_a - \epsilon_b} = \frac{\bar{g}_{ab}^{ij}}{\Delta_{ab}^{ij}}
\end{align}
are the antisymmetrized MP2 amplitudes, where $\epsilon_p$ are the orbital energies. Indices $i$ and $j$ denote occupied orbitals, while indices $a$ and $b$ denote virtual orbitals.
Note:
\begin{align}
    {g}_{ij}^{ab} =  {g}_{ji}^{ba} = {g^*}_{ab}^{ij} =   {g^*}_{ba}^{ji}
\end{align}
Before we look at spin-separated $E_{\textbf{MP2}}$, we need to realize that not all spin-combinations are non-zero. This can be understood by working out the spin-block separation of $\bar{g}_{ij}^{ab}$:
\begin{align}
    \bar{g}_{ij}^{ab} &= \bar{g}_{i_\alpha j_\alpha}^{a_\alpha b_\alpha} + \bar{g}_{i_\beta j_\beta}^{a_\beta b_\beta} + \bar{g}_{i_\alpha j_\beta}^{a_\alpha b_\beta}  + \bar{g}_{i_\beta j_\alpha}^{a_\beta b_\alpha} + \bar{g}_{i_\beta j_\alpha}^{a_\alpha b_\beta} + \bar{g}_{i_\alpha j_\beta}^{a_\beta b_\alpha} \nonumber \\
    \bar{g}_{i_\alpha j_\alpha}^{a_\alpha b_\alpha} &= g_{i_\alpha j_\alpha}^{a_\alpha b_\alpha} - g_{j_\alpha i_\alpha}^{a_\alpha b_\alpha} \quad \\
    \mathrm{ and } \quad 
    \bar{g}_{i_\alpha j_\beta}^{a_\alpha b_\beta} &=  g_{i_\alpha j_\beta}^{a_\alpha b_\beta} - 0 \quad \\ 
    (g_{j_\beta i_\alpha }^{a_\alpha b_\beta} & \text{ is zero due to violation of column spin symmetry})
\end{align}
Decomposing the energy expression into its individual spin components,
\begin{align}
    E_{\text{MP2}} &= E_{\alpha \alpha}^{\alpha \alpha}
    + E_{\beta \beta}^{\beta \beta}
    + E_{\alpha \beta}^{\alpha \beta}
    + E_{\beta \alpha}^{\beta \alpha}
    + E_{\beta \alpha}^{\alpha \beta}
    + E_{\alpha \beta}^{\beta \alpha}\label{eq:emp2-partitioned} \\
    &= \frac{1}{4} \biggl[ \bar{t}_{a_\alpha b_\alpha}^{i_\alpha j_\alpha} \bar{g}_{i_\alpha j_\alpha}^{a_\alpha b_\alpha}
    + \bar{t}_{a_\beta b_\beta}^{i_\beta j_\beta} \bar{g}_{i_\beta j_\beta}^{a_\beta b_\beta}
    +  \bar{t}_{a_\alpha b_\beta}^{i_\alpha j_\beta} \bar{g}_{i_\alpha j_\beta}^{a_\alpha b_\beta} 
    +  \bar{t}_{a_\beta b_\alpha}^{i_\beta j_\alpha} \bar{g}_{i_\beta j_\alpha}^{a_\beta b_\alpha} 
    + \bar{t}_{a_\beta b_\alpha}^{i_\alpha j_\beta} \bar{g}_{i_\alpha j_\beta}^{a_\beta b_\alpha}
    + \bar{t}_{a_\alpha b_\beta}^{i_\beta j_\alpha} \bar{g}_{i_\beta j_\alpha}^{a_\alpha b_\beta} \biggr]
    \label{eq:MP2-spin-split-expand}
\end{align}
Note that $E_{\alpha \beta}^{\alpha \beta} = E_{\beta \alpha}^{\beta \alpha}$ in \cref{eq:MP2-spin-split-expand}. This can be seen by considering the fourth term in \cref{eq:MP2-spin-split-expand}, which, upon swapping the columns of both $\bar{t}$ and $\bar{g}$ and then swapping the dummy indices, becomes identical to the third term in \cref{eq:MP2-spin-split-expand}. Similarly, it can also be seen that $E_{\beta \alpha}^{\alpha \beta}= E_{\alpha \beta}^{\beta \alpha}$. Consequently, \cref{eq:MP2-spin-split-expand} can be rewritten as
\begin{align}
    E_{\text{MP2}} 
    &= \frac{1}{4} \bar{t}_{a_\alpha b_\alpha}^{i_\alpha j_\alpha} \bar{g}_{i_\alpha j_\alpha}^{a_\alpha b_\alpha}
     + \frac{1}{4} \bar{t}_{a_\beta b_\beta}^{i_\beta j_\beta} \bar{g}_{i_\beta j_\beta}^{a_\beta b_\beta} 
     + \frac{1}{2} \bar{t}_{a_\alpha b_\beta}^{i_\alpha j_\beta} \bar{g}_{i_\alpha j_\beta}^{a_\alpha b_\beta} 
     + \frac{1}{2} \bar{t}_{a_\beta b_\alpha}^{i_\alpha j_\beta} \bar{g}_{i_\alpha j_\beta}^{a_\beta b_\alpha}.
         \label{eq:MP2-spin-split-compact}
\end{align}

Now consider the tensor network for $E_{\alpha \alpha}^{\alpha \alpha}$ only:
\begin{align}
    E_{\alpha \alpha}^{\alpha \alpha}
    &= \frac{1}{4} \bar{t}_{a_\alpha b_\alpha}^{i_\alpha j_\alpha} \bar{g}_{i_\alpha j_\alpha}^{a_\alpha b_\alpha}\\
    &= \frac{1}{4} 
       \bigl(  t_{a_\alpha b_\alpha}^{i_\alpha j_\alpha} - t_{b_\alpha a_\alpha}^{i_\alpha j_\alpha} \bigr)
       \bigl( g_{i_\alpha j_\alpha}^{a_\alpha b_\alpha} - g_{j_\alpha i_\alpha}^{a_\alpha b_\alpha} \bigr) \\
    &= \frac{1}{4} \bigl[ 
        t_{a_\alpha b_\alpha}^{i_\alpha j_\alpha} g_{i_\alpha j_\alpha}^{a_\alpha b_\alpha}
      - t_{a_\alpha b_\alpha}^{i_\alpha j_\alpha} g_{j_\alpha i_\alpha}^{a_\alpha b_\alpha}  
      - t_{b_\alpha a_\alpha}^{i_\alpha j_\alpha} g_{i_\alpha j_\alpha}^{a_\alpha b_\alpha}
      + t_{b_\alpha a_\alpha}^{i_\alpha j_\alpha} g_{j_\alpha i_\alpha}^{a_\alpha b_\alpha} 
      \bigr] .
      \label{eq:emp2-aa-fully-expanded}
\end{align}
The last term in \cref{eq:emp2-aa-fully-expanded} is equal to the first term, which can be seen by swapping the column indices of $g$ and then exchanging the dummy indices $a_\alpha$ and $b_\alpha$. An analogous relation holds between the second and third terms in \cref{eq:emp2-aa-fully-expanded}. Therefore, the expression for $E_{\alpha \alpha}$ can be written in a compact form as
\begin{align}
    E_{\alpha \alpha}^{\alpha \alpha}
    &= \frac{1}{2} \bigl[ 
        t_{a_\alpha b_\alpha}^{i_\alpha j_\alpha} g_{i_\alpha j_\alpha}^{a_\alpha b_\alpha}
      - t_{a_\alpha b_\alpha}^{i_\alpha j_\alpha} g_{j_\alpha i_\alpha}^{a_\alpha b_\alpha}  
    \bigr] \\
    &= \frac{1}{2} t_{a_\alpha b_\alpha}^{i_\alpha j_\alpha}
       \bigl( g_{i_\alpha j_\alpha}^{a_\alpha b_\alpha} - g_{j_\alpha i_\alpha}^{a_\alpha b_\alpha} \bigr)\label{eq:emp2-aaaa}.
\end{align}
Similarly, the expression for $E_{\beta \beta}^{\beta \beta}$ can be written in compact form as
\begin{align}\label{eq:emp2-bbbb}
    E_{\beta \beta}^{\beta \beta}
    &= \frac{1}{2} t_{a_\beta b_\beta}^{i_\beta j_\beta} 
       \bigl( g_{i_\beta j_\beta}^{a_\beta b_\beta} - g_{j_\beta i_\beta}^{a_\beta b_\beta} \bigr) .
\end{align}
Now consider $2E_{\alpha \beta}^{\alpha \beta}$:
\begin{align}
    2E_{\alpha \beta}^{\alpha \beta} 
    &= \frac{1}{2} \bar{t}_{a_\alpha b_\beta}^{i_\alpha j_\beta} \bar{g}_{i_\alpha j_\beta}^{a_\alpha b_\beta} \\
    &= \frac{1}{2}
       \bigl(t_{a_\alpha b_\beta}^{i_\alpha j_\beta} - t_{b_\beta a_\alpha}^{i_\alpha j_\beta}\bigr)
       \bigl(g_{i_\alpha j_\beta}^{a_\alpha b_\beta} - g_{j_\beta i_\alpha}^{a_\alpha b_\beta}\bigr).
       \label{eq:emp2-ab-fully-expanded}
\end{align}
The second term in each bracket vanishes because it violates column spin equivalence, and hence $2E_{\alpha \beta}^{\alpha \beta}$ reduces to
\begin{align}\label{eq:emp2-abab}
    2E_{\alpha \beta}^{\alpha \beta}
    &= \frac{1}{2} t_{a_\alpha b_\beta}^{i_\alpha j_\beta} g_{i_\alpha j_\beta}^{a_\alpha b_\beta}.
\end{align}
Similarly, for $2E_{\beta \alpha}^{\alpha \beta}$, the first terms in the antisymmetrization brackets vanish, and we are left with
\begin{align}\label{eq:emp2-abba}
    2E_{\beta \alpha }^{\alpha \beta}
    &= \frac{1}{2} t_{a_\alpha b_\beta}^{i_\alpha j_\beta} g_{i_\alpha j_\beta}^{a_\alpha b_\beta}.
\end{align}
Therefore, substituting the results of \cref{eq:emp2-aaaa,eq:emp2-bbbb,eq:emp2-abab,eq:emp2-abba} into \cref{eq:MP2-spin-split-compact}, we get the expression for the total unrestricted MP2 energy as:
\begin{align}\label{eq:EUMP2}
    E_{\text{UMP2}} = \frac{1}{2} t_{a_\alpha b_\alpha}^{i_\alpha j_\alpha}
       \bigl( g_{i_\alpha j_\alpha}^{a_\alpha b_\alpha} - g_{j_\alpha i_\alpha}^{a_\alpha b_\alpha} \bigr) + \frac{1}{2} t_{a_\beta b_\beta}^{i_\beta j_\beta} 
       \bigl( g_{i_\beta j_\beta}^{a_\beta b_\beta} - g_{j_\beta i_\beta}^{a_\beta b_\beta} \bigr) + t_{a_\alpha b_\beta}^{i_\alpha j_\beta} g_{i_\alpha j_\beta}^{a_\alpha b_\beta}
\end{align}

\section{Spin-restricted non-relativistic MP2 energy}
The spin-restricted expression for MP2 energy can be obtained from \cref{eq:EUMP2} by simply dropping the spin indices as
\begin{align}
    E_{\text{RMP2}} &= \frac{1}{2} t_{a b}^{i j}
       \bigl( g_{i j}^{a b} - g_{j i}^{a b} \bigr) + \frac{1}{2} t_{a b}^{i j} 
       \bigl( g_{i j}^{a b} - g_{j i}^{a b} \bigr) + t_{a b}^{i j} g_{i j}^{a b} \\
       &=t_{a b}^{i j}
       \bigl( 2g_{i j}^{a b} - g_{j i}^{a b} \bigr)\label{eq:ERMP2}
\end{align}

\section{Kramers restricted relativistic MP2 energy}

If $P$ denotes a spinor (Kramers up), then let $\bar{P}$ denote\footnote{Note: $\bar{}$ over an index denotes time reversal and $\bar{}$ over an operator denotes antisymmetrization} the time reversed version (Kramers down) of $P$ , i.e.,
\begin{align}
   |P \rangle = \mathcal{T}| \bar{P} \rangle; \quad \{  |P \rangle \} \cup \{  |\bar{P} \rangle \} = \{  |p \rangle \}
\end{align}
where $\mathcal{T}$ is the time reversal operator:
\begin{align}
    \mathcal{T} = \begin{cases}
        i\boldsymbol{\sigma_y} \hat{K}_0, &  \text{2-component}\\
        i\boldsymbol{\Sigma_y} \hat{K}_0, &   \text{4-component}  
    \end{cases}
\end{align}
and $\hat{K}_0$ is the complex conjugation operator.



Due to the symmetry of the matrix elements under time reversal, the following relationships arise for electron repulsion integrals for closed shell systems:
\begin{align}
    {g}_{PQ}^{RS} ={g^*}_{\bar{P}\bar{Q}}^{\bar{R}\bar{S}} &= {g^*}_{\bar{P}S}^{\bar{R}Q} = {g^*}_{R\bar{Q}}^{P\bar{S}}\\
    -{g^*}_{ \bar{Q}R }^{ SP } = -{g^*}_{ R \bar{Q} }^{ P S } = {g}_{ Q P }^{ \bar{S}R } = {g}_{ P Q }^{ R\bar{S} } &= -{g^*}_{ S \bar{R} }^{ \bar{Q}\bar{P} } = -{g^*}_{ \bar{R} S }^{ \bar{P} \bar{Q}} = {g^*}_{ \bar{S} \bar{P} }^{ Q \bar{R} } = {g^*}_{ \bar{P} \bar{S}  }^{ \bar{R} Q } \\
    {g^*}_{ \bar{R} \bar{S}  }^{ P Q } &= {g}_{ P Q }^{ \bar{R} \bar{S} } \\
    {g^*}_{ \bar{P} Q  }^{ R \bar{S} } &= {g}_{ P \bar{Q}  }^{ \bar{R} S }
\end{align}

Let's rewrite the MP2 energy expression, separating it into each Kramers pair permuted blocks,
\begin{align}
    E_{\text{MP2}} 
    = \frac{1}{4} \biggl[ &\bar{t}_{ A B }^{ I J } \bar{g}_{ I J }^{ A B } \\
    + &\bar{t}_{ \bar{A} B }^{ I J } \bar{g}_{ I J }^{ \bar{A} B } 
    + \bar{t}_{ A \bar{B} }^{ I J } \bar{g}_{ I J }^{ A \bar{B} } 
    + \bar{t}_{ A B }^{ \bar{I} J } \bar{g}_{ \bar{I} J }^{ A B }
    + \bar{t}_{ A B }^{ I \bar{J} } \bar{g}_{ I \bar{J} }^{ A B }\\
    + &\bar{t}_{ \bar{A} \bar{B} }^{ I J } \bar{g}_{ I J }^{ \bar{A} \bar{B} }
    + \bar{t}_{ A B }^{ \bar{I} \bar{J} } \bar{g}_{ \bar{I} \bar{J} }^{ A B } 
    + \bar{t}_{ \bar{A} B }^{ \bar{I} J } \bar{g}_{ \bar{I} J }^{ \bar{A} B } 
    + \bar{t}_{ A \bar{B} }^{ I \bar{J} } \bar{g}_{ I \bar{J} }^{ A \bar{B} }
    + \bar{t}_{ A \bar{B} }^{ \bar{I} J } \bar{g}_{ \bar{I} J }^{ A \bar{B} }
    + \bar{t}_{ \bar{A} B }^{ I \bar{J} } \bar{g}_{ I \bar{J} }^{ \bar{A} B } \\
    +& \bar{t}_{ A \bar{B} }^{ \bar{I} \bar{J} } \bar{g}_{ \bar{I} \bar{J} }^{ A \bar{B} } 
    + \bar{t}_{ \bar{A} B }^{ \bar{I} \bar{J} } \bar{g}_{ \bar{I} \bar{J} }^{ \bar{A} B }
    + \bar{t}_{ \bar{A} \bar{B} }^{ I \bar{J} } \bar{g}_{ I \bar{J} }^{ \bar{A} \bar{B} }
    + \bar{t}_{ \bar{A} \bar{B}  }^{ \bar{I} J } \bar{g}_{ \bar{I} J }^{ \bar{A} \bar{B}  } \\
    +& \bar{t}_{ \bar{A} \bar{B} }^{ \bar{I} \bar{J} } \bar{g}_{ \bar{I} \bar{J} }^{ \bar{A} \bar{B} }
    \biggr]
    \label{eq:MP2-spin-split-expand}
\end{align}

\section{Reduction of the 16-block expansion}
\label{sec:KR-MP2-reduction}

The 16-block expansion of \cref{eq:MP2-spin-split-expand} reduces to a
much smaller, structurally compact form by exploiting two properties:

\begin{enumerate}[leftmargin=*]
    \item \textbf{Time-reversal symmetry (TRS) folding.} For closed-shell
    Kramers-restricted references, summands related by full
    bar-flip $X \leftrightarrow \bar{X}$ on every index are complex
    conjugates of one another. Each pair collapses via
    $A + A^{*} = 2\,\Re A$ (or $A - A^{*} = 2i\,\Im A$), so each
    conjugate-pair becomes a single $\Re$- or $\Im$-marked summand.

    \item \textbf{Antisymmetrizer recombination.} Direct/exchange pairs
    of non-symmetric tensors produced by the trace merge back into
    single antisymmetric tensors $\bar{g}$, $\bar{t}$ over the
    Kramers-restricted index space.
\end{enumerate}

The resulting six-term Kramers-restricted MP2 energy expression
($I, J, A, B$ are Kramers-up dummy indices, $\bar{I}, \bar{A}, \dots$
are their Kramers-down partners; $\bar{g}, \bar{t}$ denote
antisymmetric tensors; $\Re[\cdot]$ marks an evaluation-time real-part
of a closed contraction) is:
\begin{align}
    E_{\text{MP2}}^{\text{KR}} =\;\;
    & -\,\Re\!\left[\, g^{A_1 A_2}_{I_1 I_2}\,
                       \bar{t}^{I_2 I_1}_{A_1 A_2}\, \right] \nonumber\\
    & +\,\Re\!\left[\, \bar{g}^{A_1 \bar{A}_1}_{I_1 \bar{I}_1}\,
                       \bar{t}^{I_1 \bar{I}_1}_{A_1 \bar{A}_1}\, \right] \nonumber\\
    & +\,\Re\!\left[\, g^{\bar{A}_1 \bar{A}_2}_{I_1 I_2}\,
                       \bar{t}^{I_1 I_2}_{\bar{A}_1 \bar{A}_2}\, \right] \nonumber\\
    & +\,2\,\Re\!\left[\, g^{A_1 \bar{A}_1}_{I_1 I_2}\,
                          t^{I_1 I_2}_{A_1 \bar{A}_1}\, \right] \nonumber\\
    & -\,2\,\Re\!\left[\, g^{A_1 \bar{A}_1}_{\bar{I}_1 \bar{I}_2}\,
                          t^{\bar{I}_2 \bar{I}_1}_{A_1 \bar{A}_1}\, \right] \nonumber\\
    & -\,2\,\Re\!\left[\, g^{A_1 A_2}_{I_1 \bar{I}_1}\,
                          \bar{t}^{I_1 \bar{I}_1}_{A_2 A_1}\, \right] .
    \label{eq:KR-MP2-six-terms}
\end{align}

The structure makes the cost breakdown explicit: one
doubly-antisymmetrized $\bar{g}\cdot\bar{t}$ term (line 2), three
singly-antisymmetrized $g\cdot\bar{t}$ terms (lines 1, 3, 6), and two
non-antisymmetrized $g\cdot t$ terms (lines 4, 5).

\section{Complex (Re/Im) split: real-arithmetic form for closed-shell
1eX2C}
\label{sec:KR-MP2-csplit}

For a closed-shell 1eX2C reference the spin-free Coulomb operator
spin-traces each per-electron transition density to its
\emph{complex scalar} Pauli component, so every tensor in
\cref{eq:KR-MP2-six-terms} is intrinsically complex (2-component);
there is no quaternion structure to exploit and a 4-component
(quaternion) decomposition would over-decompose the expression for
zero benefit.

The \emph{minimal} real-arithmetic form is therefore the 2-component
\emph{complex split}: rewrite each complex tensor as two real tensors
\begin{align}
    g \;&=\; \Re g \;+\; i\,\Im g,
    \qquad
    g^{*} \;=\; \Re g \;-\; i\,\Im g, \\
    t \;&=\; \Re t \;+\; i\,\Im t,
    \qquad
    t^{*} \;=\; \Re t \;-\; i\,\Im t,
\end{align}
and propagate complex arithmetic symbolically. Each
$\Re[\,g\cdot t\,]$ summand of \cref{eq:KR-MP2-six-terms} expands as
\begin{equation}
    \Re\!\left[\, g \cdot t \,\right] \;=\;
        \Re g \cdot \Re t \;-\; \Im g \cdot \Im t,
    \label{eq:csplit-Re-product}
\end{equation}
which is a difference of two products of \emph{real} tensors. The full
six-term Kramers-restricted MP2 energy expression in the
complex-split form (with $\Re g$, $\Im g$, $\Re t$, $\Im t$ all real)
is:
\begin{align}
    E_{\text{MP2}}^{\text{KR,c-split}} =\;\;
    & -\,\Re\!\left[\,
        \Re g^{A_1 A_2}_{I_1 I_2}\,
        \bar{\Re t}^{I_2 I_1}_{A_1 A_2}
      - \Im g^{A_1 A_2}_{I_1 I_2}\,
        \bar{\Im t}^{I_2 I_1}_{A_1 A_2}
      \,\right] \nonumber\\
    & +\,\Re\!\left[\,
        \bar{\Re g}^{A_1 \bar{A}_1}_{I_1 \bar{I}_1}\,
        \bar{\Re t}^{I_1 \bar{I}_1}_{A_1 \bar{A}_1}
      - \bar{\Im g}^{A_1 \bar{A}_1}_{I_1 \bar{I}_1}\,
        \bar{\Im t}^{I_1 \bar{I}_1}_{A_1 \bar{A}_1}
      \,\right] \nonumber\\
    & +\,\Re\!\left[\,
        \Re g^{\bar{A}_1 \bar{A}_2}_{I_1 I_2}\,
        \bar{\Re t}^{I_1 I_2}_{\bar{A}_1 \bar{A}_2}
      - \Im g^{\bar{A}_1 \bar{A}_2}_{I_1 I_2}\,
        \bar{\Im t}^{I_1 I_2}_{\bar{A}_1 \bar{A}_2}
      \,\right] \nonumber\\
    & +\,2\,\Re\!\left[\,
        \Re g^{A_1 \bar{A}_1}_{I_1 I_2}\,
        \Re t^{I_1 I_2}_{A_1 \bar{A}_1}
      - \Im g^{A_1 \bar{A}_1}_{I_1 I_2}\,
        \Im t^{I_1 I_2}_{A_1 \bar{A}_1}
      \,\right] \nonumber\\
    & -\,2\,\Re\!\left[\,
        \Re g^{A_1 \bar{A}_1}_{\bar{I}_1 \bar{I}_2}\,
        \Re t^{\bar{I}_2 \bar{I}_1}_{A_1 \bar{A}_1}
      - \Im g^{A_1 \bar{A}_1}_{\bar{I}_1 \bar{I}_2}\,
        \Im t^{\bar{I}_2 \bar{I}_1}_{A_1 \bar{A}_1}
      \,\right] \nonumber\\
    & -\,2\,\Re\!\left[\,
        \Re g^{A_1 A_2}_{I_1 \bar{I}_1}\,
        \bar{\Re t}^{I_1 \bar{I}_1}_{A_2 A_1}
      - \Im g^{A_1 A_2}_{I_1 \bar{I}_1}\,
        \bar{\Im t}^{I_1 \bar{I}_1}_{A_2 A_1}
      \,\right] .
    \label{eq:KR-MP2-csplit}
\end{align}

The outer $\Re[\cdot]$ wrappers are now structurally redundant ---
every inner product is a product of real tensors and therefore real
--- but they are retained as evaluation-dispatch markers so the inner
contractions can be performed as real (DGEMM) rather than complex
(ZGEMM) products. This is the cost win of the complex split versus
\cref{eq:KR-MP2-six-terms}.

\end{document}
